\documentclass{article}
\usepackage{listings}
\usepackage{amsmath,amsthm}
\usepackage{amssymb}
\usepackage{hyperref}
\usepackage{xurl}


\title{Group projects Programming in Python II}
\author{}
\date{}

\begin{document}
\maketitle

\textbf{
   The project consists of two parts: the project itself and a short quiz, which will be available later.
   This quiz will be different for each team,
   and it will be based on the submitted project. The quiz will be made available shortly after the submission of the project and it will be based on the
   submission. The quiz will consists of questions on the submitte code: you will be asked to explain some of your code or indicate ho
   w to change it.
   The quiz will be timed (max. 15 minutes), it will be available on K2, and you will be able to do it online, within a few days after the submission of your project. }
   
   \textbf{You are allowed to use anything for your project, including generative AI. However, you have to make sure that you are able to explain your code and answer questions on it
   during the quiz. }

  \textbf{ Deadline for project: 21/12/2025, 23:59,  please submit your *.py or *.ipynb files using the designated folder on the K2 site of the course.  Quiz available
   between 23/12/2025-28/12/2025.
}



\section*{Project 1: Group 2 (Callum Murphy,David Tomasso),
Group 13 (Ahimin Paul Martin AMETHIER,Benjamin  AMSTUTZ, Roy-Calvin BARBEN)
Group 14 (Ekaterina BOGDANOVA, Polina DRONOVA, Jocelyne GABELA SANTANDER)
}
The task is to predict stock prices based on past stock prices. The data is stored in the file \lstinline$tech_companies_stock_prices.csv$.

The fields are as follows:
\begin{itemize} \item \lstinline$Date$ -- the date of the transaction 
\item \lstinline$Stock_name$ -- the name of the company 
\item \lstinline$Open$, \lstinline$Close$ -- opening and closing prices on the transaction date 
\item \lstinline$Open_i$, \lstinline$Close_i$ -- for $i=1,2,3,4,5$, these represent the opening and closing prices on \lstinline$Date$-$i$ 
\end{itemize}

For each company, predict the share price based on its past share prices over the previous $k$ dates. Specifically: 
\begin{itemize} 
\item Try to predict both the opening and closing prices using various choices of past opening and closing prices. 
\item Explore using past share prices of other companies to predict the share price of a given company. 
\end{itemize}

Use various models (such as linear regression, SVM, decision trees, and neural networks) to predict share prices. Compare the models in terms of prediction accuracy to determine which model performs best.

Finally, select one or two independent variables and use them to build models for predicting share prices. Identify which variables provide the best model for predicting share prices.



\section*{Project 2:  Group 3 (Kim Vast, Constance Thouvenot, Célia Pinto Moreira), Group 4 (Yali Xiao,
Boxi Li),  Group 10 (Léa SOKALSKI, Chanelle STARK,Yasmine ZEMZEM), Group 15 (Chloe BUNDA, Medha	CHAUHAN,Shreyas	BHUTRA),
Group 16 (I-An CHEN, Himanshu CHORARIA, Gabriele COLI)
}
The task is to predict stock prices based on past stock prices. The data is stored in the file \lstinline$stock_data_for_hw1.csv$.

The fields are as follows: \begin{itemize} \item \lstinline$Date$ -- date of the transaction \item \lstinline$Company$ -- name of the company \item \lstinline$Open$, \lstinline$Close$, \lstinline$High$, \lstinline$Low$ -- opening, closing, highest, and lowest prices on the transaction date \item \lstinline$Open_i$, \lstinline$Close_i$, \lstinline$High_i$, \lstinline$Low_i$ -- for $i=1,2,3,4,5$, these represent the opening, closing, highest, and lowest prices on \lstinline$Date$-$i$ \end{itemize}

For each company, predict the share price based on its past share prices over the previous $k$ dates. Specifically: \begin{itemize} \item Try to predict opening, closing, high, and low prices separately, using different combinations of past opening, closing, high, and low prices. \item Investigate using past share prices of other companies to predict the share price of a given company. \end{itemize}

Use various models (e.g., linear regression, SVM, decision trees, neural networks) for predicting share prices. Compare the models in terms of prediction accuracy to determine which one performs best.

Finally, select one or two independent variables to build models for predicting share prices. Identify which variables yield the best predictive model.







\section*{Project 3:  Group 5 (Jinjun Tang, Diwen Shi), Group 6 (Vengateshwaran NADARAJAN, ZOIROV Bahodir, SHIJI SUNIL Samudra), Group 9 (Carl Shiminyi Neba MOHMBAMEMBIE, Shravan SIVAKUMAR,Tkhui Lin ZYONH),
Group 17 (Pierre-Henri Jean  André COUFFON DE TREVROS, Pauline D'ARRAS, Samyakranjan DAS),
Group 18 (Côme DE LANGSDORFF, Mahault DE TARRAGON, Louis DERVEAUX)
}
 Predict stock prices based on past stock prices, ESG Score, and net income. The CSV file \lstinline$StockESG0.csv$ contains the following columns: \begin{itemize} \item \lstinline$Date$ -- date of transaction \item \lstinline$Company$ -- name of the company \item \lstinline$Open$, \lstinline$Close$ -- opening and closing prices on the transaction date \item \lstinline$Open_i$, \lstinline$Close_i$ -- for $i=1,2,3,4,5$, these represent the opening and closing prices on \lstinline$Date$-$i$ \item \lstinline$Net income$ -- net income of the company \item \lstinline$ESG Score$ -- overall ESG score \item \lstinline$ESG Controversies Score$, \lstinline$Environmental Controversies Count$, \lstinline$Human Rights Controversies$, \lstinline$Business Ethics Controversies$, \lstinline$Wages Working Condition Controversies Count$, \lstinline$Total Renewable Energy To Energy Use in million$, \lstinline$Estimated CO2 Equivalents Emission Total$ -- various ESG-related scores and metrics \end{itemize}

Use past opening or closing prices, net income, and various ESG scores to predict the current opening (or closing) price.

Apply different models, such as linear regression, SVM, decision trees, and neural networks, for predicting the share price. Compare the models in terms of prediction accuracy to determine which performs best.

Finally, select one or two independent variables to build predictive models for share prices. Identify which variables yield the most accurate model. 



 \section*{Project 4:  Group 7 (Guoyu ZHU, Dayu LI), Group 8 (Jingyi LI, Yichen ZHOU), Group 11 (Uddesh Uddhav	TANDEL,Anjali	MANOJ, Xiao WANG),
  Group 19 (Valentin	DIGUET, Thibault DINDAULT, Kenneth DJINADOU),
  Group 20 (Marc DORADOUX, Ilyas FADEL, Yuxuan	FAN), Group 21 (Aymeric	GORISSE, Xin HU, Haodan	HUANG)
 }
The task is to predict the ESG score based on various other variables, such as: \begin{itemize} \item \lstinline$ESG Controversies Score$ \item \lstinline$Environmental Controversies Count$ \item \lstinline$Human Rights Controversies$ \item \lstinline$Business Ethics Controversies$ \item \lstinline$Wages Working Condition Controversies Count$ \item \lstinline$Total Renewable Energy To Energy Use in million$ \item \lstinline$Estimated CO2 Equivalents Emission Total$ \end{itemize}

The data is stored in the file \lstinline$ESGAllCompanies1.csv$.

Use various models (e.g., linear regression, SVM, decision trees, neural networks) to predict the ESG score. Compare the performance of the models in terms of prediction accuracy. Which model performs the best?

Select one or two independent variables and build models to predict the ESG score. Identify which variables provide the best predictive model for the ESG score. 

\section*{Project 5:  Group 12 (Damien LOUSTEAU,Baihass SABET AGHA),  Group 22 (Andreas	JESPERSEN, Amelie JIN), Group 1 (Maria LAMEIRAS FRANCO SOARES PINTO, Gregorio MARCHESE, Filip LOUER)}
   The task is to reverse engineer a portfolio management algorithm. 
   The historial  data is stored in the file \lstinline$tech_companies_stock_prices.csv$. The fields are as follows:
\begin{itemize} \item \lstinline$Date$ -- the date of the transaction 
\item \lstinline$Stock_name$ -- the name of the company 
\item \lstinline$Open$, \lstinline$Close$ -- opening and closing prices on the transaction date 
\item \lstinline$Open_i$, \lstinline$Close_i$ -- for $i=1,2,3,4,5$, these represent the opening and closing prices on \lstinline$Date$-$i$ 
\end{itemize}

The historical performance of the portfolio management algorithm can be found in \lstinline$test_portfolio.csv$. 
The rows of this file are indexed by dates of transactions, the columns correspond to various stocks. In addition, there is a 
column \lstinline$cash$ and a column \lstinline$value$. The column \lstinline$cash$ represents the cash added to  (if positive) or spent from (if negative) the portfolio at
the beginning of the given day. The columns \lstinline$value$ represents the value of the portfolio at the end of the given day.
The values in columns labeled by stocks correspond to the number of stocks bought (if positive) or sold (if negative) in the 
begining of the day, using the opening prices. For computing the value the cash in the portfolio and the values of the stocks in the
portfolio for closing prices are added. 
It is assumed that at the beginning of the day the proceeds from selling stocks at the opening prices and the cash reserve 
can be used to buy stock at the opening prices. 

For example, the line
\begin{lstlisting}
	,    TSLA, AAPL, AMZN, MSFT, NVDA, META, GOOGL,cash,               value
  2023-10-10,0.0,  4.0,  -2.0, -2.0, 0.0,  0.0,  1.0,   68.65997314452716, 303305.98236083984
\end{lstlisting}
means that on \lstinline$2023-10-10$, at the beginning of the day, $4$ more shares of 'AAPL' were bought, $2$ shares of 'AMZN' and 'MSFT' were
sold, and $1$ more share of 'GOOGL' was bought and $68.65997314452716$ was added to the cash reserves. At the end of the day,
the value of the portfolio was ' 303305.98236083984', where the closing prices were used for computing the value.
Note that $4$ is not the number of 'AAPL' share in the portfolio, but the number of 'AAPL' shares which were bought in addition to
the 'AAPL' shares already in the portfolio. Likewise, '-2.0' is the number of 'AMZN' shares sold from the number of 'AMZN' shares in
the portfolio. 


First, use the historical data from \lstinline$tech_companies_stock_prices.csv$ and \lstinline$test_portfolio.csv$
to reconstruct the portfolio (number of stocks, cash) at the end of each trading date. Second, various models 
(e.g., linear regression, SVM, decision trees, neural networks) to predict 
\begin{itemize}
\item
    the value of the portfolio at the end of each day (except the first trading day) 
     as a function of the portfolio at the end of the previous day and the opening prices of the current day.
\item
    the trading actions (buying/selling stocks) as a function of the trading actions during the previous $k$ days ($k=1,2,3$).
\end{itemize}

Keep at least 20\% of the data for testing.

Try to propose a better trading strategy than the one implemented in \lstinline$test_portfolio.csv$, and compare its performance
on the test data, which was not used for reverse engineering of the trading strategy.

\end{document}
The task is to predict the logarithm of house prices (variable \emph{lprice})d using the following variables: \begin{itemize} \item logarithm of the lot size (variable \emph{llotsize}) \item logarithm of the surface area of the house (variable \emph{lsqrft}) \item number of bedrooms (variable \emph{bdrms}) \item surface area of the house (variable \emph{sqrft}) \item lot size (variable \emph{lotsize}) \end{itemize}

The data is stored in the file \lstinline$hprice1_merged.csv$.

Use various models (such as linear regression, SVM, decision trees, and neural networks) to predict the logarithm of the house price. Compare the performance of the models in terms of prediction accuracy. Which model performs the best?

Select one or two independent variables and build models to predict the logarithm of the price. Identify which variables provide the best model for predicting the logarithm of the price.

Repeat the exercise by replacing the logarithm of the price with the price itself (variable \emph{price}).


\end{document}
